\section{Conclusion}
\label{sec:conclusion}

In this work, we present unified multimodal generation as a formulation for computer vision, analogous to the role of GPT-style generative modeling in NLP.
We present the SenseNova-Vision Corpus by casting heterogeneous computer vision annotations into text, image, and mixed text-and-image generation targets, and train SenseNova-Vision within the same formulation.
This conversion enables large-scale training of a single UMM across structured visual understanding, dense geometric prediction, segmentation, and multi-view visual geometry without task-specific heads.
The resulting model matches leading task-specialized systems and supports language-defined task variants, suggesting a scalable route for absorbing computer vision supervision into general-purpose foundation models.
The generative modeling of diverse 2D and 3D perception tasks may further cultivate implicit spatial understanding, connecting this paradigm naturally to the rapidly emerging frontiers of physical intelligence.

This formulation opens several directions for future work.
First, stronger in-context learning could further reduce the boundaries between task domains, enabling new visual tasks to be specified by examples, prompts, or mixed demonstrations beyond predefined datasets.
Second, extending unified multimodal generation from images to video would bring temporal dynamics and web-scale video supervision into foundation-model training.
Finally, scaling the corpus and model capacity, together with deeper integration with the strongest language models, may allow general-purpose foundation models to absorb richer visual and spatial knowledge from computer vision, moving toward world models that can perceive, reason about, and interact with the physical world.
